\begin{table*}[t]
\begin{adjustwidth}{-\extralength}{0cm}
\centering
\footnotesize
\setlength{\tabcolsep}{3pt}
\caption{Retrieval-filter detection quality.\label{tab:detection}}
\begin{tabular}{lrrrrrrrrr}
\toprule
\textbf{Attack} & \textbf{SQuAD AUROC} & \textbf{AUPRC} & \textbf{TPR} & \textbf{HotpotQA AUROC} & \textbf{AUPRC} & \textbf{TPR} & \textbf{Toy corpus AUROC} & \textbf{AUPRC} & \textbf{TPR}\\
\midrule
PoisonedRAG & 0.99 & 0.99 & 0.99 & 0.97 & 0.93 & 0.96 & 0.99 & 0.99 & 1.00\\
Trigger backdoor & 1.00 & 1.00 & 1.00 & 1.00 & 1.00 & 1.00 & 1.00 & 1.00 & 1.00\\
Prompt injection & 1.00 & 0.99 & 1.00 & 1.00 & 0.97 & 1.00 & 0.99 & 0.99 & 1.00\\
Context flooding & 0.99 & 1.00 & 0.99 & 0.98 & 0.98 & 0.97 & 0.99 & 1.00 & 1.00\\
Adaptive poisoning & 0.47 & 0.56 & 0.00 & 0.65 & 0.47 & 0.02 & 0.31 & 0.65 & 0.00\\
\bottomrule
\end{tabular}
\noindent\footnotesize Measured in the filter-only configuration so that Layer 1 does not remove the candidates the filter is being scored on. AUROC and average precision (AUPRC) are threshold-free and computed over every retrieval candidate; TPR is measured at the deployed operating point. Poisoned candidates are a small minority of the candidate set, so AUPRC is the more informative summary. False-positive rate on unattacked traffic was 0.01 (SQuAD), 0.01 (HotpotQA), 0.00 (Toy corpus).
\end{adjustwidth}
\end{table*}
